% Vietnamese translation for OI Public Library
% Source-File: IOI中国国家候选队论文/2005/唐文斌/唐文斌.ppt
\documentclass[11pt]{article}
\usepackage{oipl-vn}

\newcommand{\Oh}{\mathcal{O}}

\title{Bản trình chiếu: Tư duy nghịch hướng trong giải bài}
\author{Tang Wenbin}
\date{2005}

\begin{document}
\maketitle

\origtitle{Khó theo chiều thuận thì xét chiều ngược: bản trình chiếu}
\sourcepath{IOI China candidate team papers / 2005 / Tang Wenbin / Tang Wenbin.ppt}

\section{Mạch chính của slide}

Bản trình chiếu đi cùng bài viết \emph{Khó theo chiều thuận thì xét chiều
ngược}. Phần chữ trích xuất được từ PowerPoint cũ giữ lại ba ví dụ chính:
\textit{Dinner Is Ready}, \textit{Greedy Path} và \textit{Building Towers}.
Các công thức còn đọc được gồm hệ
\[
  X_1+X_2+\cdots+X_n=M,
\]
các ràng buộc \(\mathit{Mini},\mathit{Maxi}\), biểu thức tổ hợp
\(\binom{M+n-1}{n-1}\), phép bù trừ trên các tập \(\overline{S_i}\), hàm
\(o(t)\) cho chu trình tỉ lệ, và quy hoạch động \(F[i,j,k]\) cho bài xây
tháp. Bản dịch này dùng bài viết đã dịch làm neo nội dung, nhưng giữ giọng
của slide: nhấn vào cách đổi hướng suy nghĩ, không lặp lại toàn bộ phụ lục đề
bài hay chương trình kèm theo.

Thông điệp trung tâm là: khi mô hình thuận dẫn tới quá nhiều trường hợp hoặc
quá nhiều trạng thái, hãy thử nhìn vào phần bù, vào điều kiện của đáp án, hoặc
vào trạng thái cuối rồi đi ngược lại. ``Nghịch hướng'' ở đây không phải một
thuật toán cố định; nó là thao tác đổi đối tượng cần xét để bài toán lộ ra
cấu trúc dễ tính hơn.

\section{Đếm phần bù trong \textit{Dinner Is Ready}}

Ví dụ đầu tiên là một bài đếm nghiệm nguyên. Có \(n\) đứa trẻ và \(M\) khúc
xương; đứa trẻ thứ \(i\) nhận \(X_i\) khúc, với
\[
  \mathit{Min}_i\le X_i\le \mathit{Max}_i.
\]
Nếu chỉ có giới hạn dưới, phép đổi biến
\[
  Y_i=X_i-\mathit{Min}_i
\]
đưa bài toán về đếm nghiệm không âm, và số nghiệm của
\[
  Y_1+\cdots+Y_n=M-\sum_i\mathit{Min}_i
\]
có thể tính bằng công thức tổ hợp. Khó khăn còn lại là các giới hạn trên.
Đếm trực tiếp các nghiệm vừa đủ mọi giới hạn trên không thuận tiện, nhất là
khi \(n\) tuy nhỏ nhưng mỗi biến có biên riêng.

Slide vì vậy chuyển sang mặt ngược: đếm các nghiệm vi phạm ít nhất một giới
hạn trên. Gọi \(S\) là tập nghiệm chỉ thỏa giới hạn dưới, \(S_i\) là tập nghiệm
trong \(S\) thỏa \(X_i\le \mathit{Max}_i\), và \(\overline{S_i}\) là tập nghiệm
với \(X_i>\mathit{Max}_i\). Đáp án cần tìm là
\[
  |S_1\cap S_2\cap\cdots\cap S_n|.
\]
Thay vì đếm giao này, ta dùng nguyên lý bù trừ:
\[
\begin{aligned}
&|S_1\cap\cdots\cap S_n|\\
&=|S|-\sum_i|\overline{S_i}|
  +\sum_{i<j}|\overline{S_i}\cap\overline{S_j}|-\cdots .
\end{aligned}
\]
Mỗi giao của các tập bù lại trở thành một hệ nghiệm không âm sau khi trừ thêm
phần vượt quá \(\mathit{Max}_i\). Vì \(n\le 8\), việc liệt kê các tập con của
biến là rất nhỏ; độ phức tạp được slide ghi là \(\Oh(2^n(n+M))\).

\section{Giả sử đã biết đáp án trong \textit{Greedy Path}}

Ví dụ thứ hai là bài tìm chu trình có tỉ số lớn nhất:
\[
  F(C)=\frac{\sum_{e\in C} C_e}{\sum_{e\in C} T_e}.
\]
Nếu đi theo chiều thuận, ta phải so sánh trực tiếp các chu trình, nhưng mục
tiêu là một phân số chứ không phải tổng trọng số. Slide đổi câu hỏi: nếu giá
trị tối ưu là \(t^*\), điều gì phải đúng?

Với một số thực \(t\), đặt trọng số mới cho mỗi cạnh là
\[
  W_e=C_e-tT_e
\]
và định nghĩa
\[
  o(t)=\max_C\sum_{e\in C} W_e.
\]
Nếu \(t<t^*\), tồn tại chu trình có tổng \(C_e-tT_e\) dương; nếu \(t>t^*\),
mọi chu trình đều cho tổng âm; tại \(t=t^*\), chu trình tối ưu cho giá trị
không. Như vậy bài toán tỉ lệ được chuyển thành bài toán kiểm tra dấu của
\(o(t)\).

Từ đó có thuật toán tìm kiếm tham số. Chọn khoảng ban đầu \((t_l,t_h)\) chứa
đáp án, lấy trung điểm \(t\), đổi trọng số cạnh thành \(C_e-tT_e\), rồi tìm
chu trình có tổng lớn nhất. Bản slide ghi bước này bằng Floyd và độ phức tạp
dạng \(\Oh(\log K\cdot N^3)\), trong đó \(K\) biểu thị độ chính xác hoặc số
lần chia đôi. Điểm quan trọng không nằm ở Floyd, mà ở việc ``giả sử đã biết
đáp án'' để biến một bài tối ưu phân số thành một loạt bài kiểm tra dấu.

\section{Đảo chiều quy hoạch trong \textit{Building Towers}}

Ví dụ cuối là bài đếm tháp. Mỗi tháp có \(H\) tầng, tầng đáy có \(M\) khối,
và mỗi tầng tiếp theo hơn hoặc kém tầng dưới đúng \(1\) khối. Tổng số khối
được dùng không vượt quá \(N\), với các giới hạn trong slide
\[
  N\le 32767,\qquad H\le 60,\qquad M\le 10.
\]
Hướng thuận dễ nghĩ là quy hoạch động theo tầng, độ cao hiện tại và số khối
đã dùng. Slide ghi công thức dạng
\[
  F[i,j,k]=F[i+1,j+1,k-i]+F[i-1,j+1,k-i],
\]
cùng nhận xét số trạng thái có thể tới khoảng \(60\cdot70\cdot2400\), tức
hàng chục triệu. Bộ nhớ còn khó hơn vì bài cần trả lời tháp thứ \(K\) theo
thứ tự từ điển, nên không thể chỉ cuốn chiếu bảng.

Cách nghịch hướng là tìm kiếm từ trên xuống có ghi nhớ. Trạng thái được nhìn
như bộ ba \((N,H,M)\): còn nhiều nhất \(N\) khối, còn \(H\) tầng, và tầng hiện
tại có \(M\) khối. Hàm \(F(N,H,M)\) trả về số tháp hợp lệ từ trạng thái này.
Tìm kiếm chỉ mở các nhánh thật sự có thể xuất hiện trong truy vấn, và trước
khi đi tiếp có thể cắt bỏ trạng thái vô lý bằng các kiểm tra đơn giản:
\begin{itemize}
  \item nếu số khối còn lại nhỏ hơn nhu cầu tối thiểu thì đáp án bằng \(0\);
  \item nếu số khối còn lại vượt quá nhu cầu tối đa, có thể ép nó về nhu cầu
  tối đa vì phần dư không làm thay đổi số tháp;
  \item nếu mọi dãy lên/xuống còn lại đều hợp lệ, trả ngay \(2^{H-1}\).
\end{itemize}

Ghi nhớ hóa biến tìm kiếm thành một dạng quy hoạch động, nhưng thứ tự sinh
trạng thái đã đảo lại. Các trạng thái gần gốc ít bị lặp, còn những trạng thái
sâu dễ lặp được lưu trong bảng băm. Bài viết gốc còn nêu kinh nghiệm chỉ lưu
những trạng thái có giá trị không quá một ngưỡng để giảm bộ nhớ tra cứu; slide
chỉ giữ ý chính là số trạng thái thực tế nhỏ hơn rất nhiều so với bảng thuận.

\section{Kết luận}

Ba ví dụ trong slide tương ứng với ba kiểu đổi hướng khác nhau. Với
\textit{Dinner Is Ready}, ta không đếm trực tiếp tập cần tìm mà đếm các tập
vi phạm. Với \textit{Greedy Path}, ta không cố dựng ngay chu trình tối ưu mà
hỏi giá trị tối ưu sẽ làm bài kiểm tra nào đổi dấu. Với \textit{Building
Towers}, ta không điền bảng từ đáy lên mà tìm từ trạng thái yêu cầu xuống các
trạng thái con cần thiết.

Điểm chung là cách nhìn ngược làm giảm phần dư thừa trong phân tích. Nó không
thay thế tổ hợp, tìm kiếm tham số hay quy hoạch động; nó giúp chọn đúng dạng
cho các công cụ ấy.

\end{document}
